%%%%%%%%%%%%%%%%%%%%%%%%%%%%%%%%%%%%%%%%%%%%%%%%%%%%%%%%%%%%%%%%%%%%%%%%%%%%%%%
% BEATRIX: Antworten auf Fasnacht-Fragebogen (Version 3)
% Mit EBF-Kapitel-Integration
%%%%%%%%%%%%%%%%%%%%%%%%%%%%%%%%%%%%%%%%%%%%%%%%%%%%%%%%%%%%%%%%%%%%%%%%%%%%%%%
\documentclass[11pt,a4paper]{article}

% Packages
\usepackage[utf8]{inputenc}
\usepackage[T1]{fontenc}
\usepackage[ngerman]{babel}
\usepackage{geometry}
\usepackage{graphicx}
\usepackage{booktabs}
\usepackage{array}
\usepackage{longtable}
\usepackage{xcolor}
\usepackage{tcolorbox}
\usepackage{enumitem}
\usepackage{hyperref}
\usepackage{fancyhdr}
\usepackage{lastpage}
\usepackage{amsmath}
\usepackage{amssymb}

% Page geometry
\geometry{left=2.5cm, right=2.5cm, top=2.5cm, bottom=2.5cm}

% Colors
\definecolor{innogreen}{RGB}{0, 102, 51}
\definecolor{innoblue}{RGB}{0, 51, 102}
\definecolor{innogray}{RGB}{128, 128, 128}
\definecolor{ebfpurple}{RGB}{102, 51, 153}
\definecolor{ebforange}{RGB}{204, 102, 0}

% tcolorbox styles
\tcbuselibrary{skins,breakable}

\newtcolorbox{keyinsight}{
    colback=innogreen!5!white,
    colframe=innogreen!75!black,
    title={\textbf{Kernaussage}},
    fonttitle=\bfseries,
    breakable
}

\newtcolorbox{summarybox}{
    colback=gray!10!white,
    colframe=gray!50!black,
    title={\textbf{Zusammenfassung}},
    fonttitle=\bfseries,
    breakable
}

\newtcolorbox{limitbox}{
    colback=orange!5!white,
    colframe=orange!75!black,
    title={\textbf{Ehrliche Einschränkung}},
    fonttitle=\bfseries,
    breakable
}

\newtcolorbox{ebfbox}{
    colback=ebfpurple!5!white,
    colframe=ebfpurple!75!black,
    title={\textbf{EBF-Fundierung}},
    fonttitle=\bfseries,
    breakable
}

\newtcolorbox{corebox}{
    colback=innoblue!5!white,
    colframe=innoblue!75!black,
    title={\textbf{10C CORE Framework}},
    fonttitle=\bfseries,
    breakable
}

\newtcolorbox{formulabox}{
    colback=ebforange!5!white,
    colframe=ebforange!75!black,
    fonttitle=\bfseries,
    breakable
}

% Header and footer
\pagestyle{fancy}
\fancyhf{}
\fancyhead[L]{\textcolor{innogray}{\small BEATRIX – Innosuisse Innovation Cheque}}
\fancyhead[R]{\textcolor{innogray}{\small Version 3.0 (EBF-Integration)}}
\fancyfoot[C]{\textcolor{innogray}{\small Seite \thepage\ von \pageref{LastPage}}}
\renewcommand{\headrulewidth}{0.4pt}
\renewcommand{\footrulewidth}{0.4pt}

% Hyperref setup
\hypersetup{
    colorlinks=true,
    linkcolor=innoblue,
    urlcolor=innoblue,
    citecolor=innoblue
}

%%%%%%%%%%%%%%%%%%%%%%%%%%%%%%%%%%%%%%%%%%%%%%%%%%%%%%%%%%%%%%%%%%%%%%%%%%%%%%%
\begin{document}

%%%%%%%%%%%%%%%%%%%%%%%%%%%%%%%%%%%%%%%%%%%%%%%%%%%%%%%%%%%%%%%%%%%%%%%%%%%%%%%
% TITLE PAGE
%%%%%%%%%%%%%%%%%%%%%%%%%%%%%%%%%%%%%%%%%%%%%%%%%%%%%%%%%%%%%%%%%%%%%%%%%%%%%%%

\begin{titlepage}
\centering
\vspace*{2cm}

{\Huge\bfseries\textcolor{innoblue}{BEATRIX}}\\[0.5cm]
{\Large Behavioral Economics AI Technology for\\Research and Implementation eXcellence}\\[2cm]

{\LARGE\bfseries Antworten auf Fasnacht-Fragebogen}\\[0.5cm]
{\large Version 3.0 – Mit EBF-Kapitel-Integration}\\[2cm]

\begin{tcolorbox}[colback=innoblue!10!white, colframe=innoblue!75!black, width=12cm]
\centering
\textbf{Für:} Innosuisse Innovation Cheque (CHF 15'000)\\[0.3cm]
\textbf{Einreichung bei:} SSBM (Social Sciences \& Business Management)\\[0.3cm]
\textbf{Wirtschaftspartner:} FehrAdvice Partners AG\\[0.3cm]
\textbf{Forschungspartner:} Universität Zürich\\[0.3cm]
\textbf{Datum:} 18. Januar 2026
\end{tcolorbox}

\vfill

\begin{tcolorbox}[colback=ebfpurple!10!white, colframe=ebfpurple!75!black, width=14cm]
\centering
\textbf{Wissenschaftliche Fundierung (EBF)}\\[0.3cm]
\begin{tabular}{ll}
21 Kapitel & Formales Framework \\
203+ Axiome & Dokumentierte Grundlagen \\
8 Ψ-Dimensionen & Kontextmodellierung \\
15+ γ-Parameter & Wechselwirkungen \\
\end{tabular}
\end{tcolorbox}

\vspace{1cm}
{\small\textit{Die Technologie ermöglicht die Innovation,\\aber die Transformation der Beratungspraxis ist das Produkt.}}

\end{titlepage}

%%%%%%%%%%%%%%%%%%%%%%%%%%%%%%%%%%%%%%%%%%%%%%%%%%%%%%%%%%%%%%%%%%%%%%%%%%%%%%%
% TABLE OF CONTENTS
%%%%%%%%%%%%%%%%%%%%%%%%%%%%%%%%%%%%%%%%%%%%%%%%%%%%%%%%%%%%%%%%%%%%%%%%%%%%%%%

\tableofcontents
\newpage

%%%%%%%%%%%%%%%%%%%%%%%%%%%%%%%%%%%%%%%%%%%%%%%%%%%%%%%%%%%%%%%%%%%%%%%%%%%%%%%
% STYLE GUIDE
%%%%%%%%%%%%%%%%%%%%%%%%%%%%%%%%%%%%%%%%%%%%%%%%%%%%%%%%%%%%%%%%%%%%%%%%%%%%%%%

\section*{Style-Guide Prinzipien}
\addcontentsline{toc}{section}{Style-Guide Prinzipien}

\subsection*{SSBM-Wording (Management-Fokus)}

\begin{center}
\begin{tabular}{p{6cm}p{6cm}}
\toprule
\textbf{❌ Vermeiden} & \textbf{✓ Nutzen} \\
\midrule
``Wir entwickeln einen Algorithmus...'' & ``Wir entwickeln ein Entscheidungssystem...'' \\
``KI-Architektur'' & ``Skalierbarkeit und Prozess-Sicherheit'' \\
``Software Engineering'' & ``Technologische Umsetzung'' \\
``Feature Set'' & ``Anwendungsfälle'' \\
``User Interface'' & ``Zusammenarbeit Mensch-System'' \\
\bottomrule
\end{tabular}
\end{center}

\subsection*{Fehr-Safe Language (Reputationsschutz)}

\begin{center}
\begin{tabular}{p{6cm}p{6cm}}
\toprule
\textbf{❌ Vermeiden} & \textbf{✓ Nutzen} \\
\midrule
``Revolutioniert'' & ``Weiterentwickelt systematisch'' \\
``Bahnbrechend'' & ``Vielversprechend'' \\
``Garantiert bessere Ergebnisse'' & ``Zeigt positive Resultate in Pilotprojekten'' \\
\bottomrule
\end{tabular}
\end{center}

\subsection*{Was ist neu in V3?}

Diese Version integriert die \textbf{wissenschaftliche Fundierung} aus dem Evidence-Based Framework (EBF):

\begin{center}
\begin{tabular}{lll}
\toprule
\textbf{Frage} & \textbf{Thema} & \textbf{EBF-Hauptkapitel} \\
\midrule
a) & Innovation Potential & Kap. 5, 9, 10, 11 \\
b) & Market Relevance & Kap. 12, 13, 17 \\
c) & Feasibility \& Risk & Kap. 10, 17 \\
d) & Research Partner & Kap. 11, 12, 17 \\
e) & Impact Switzerland & Kap. 1, 17 \\
\bottomrule
\end{tabular}
\end{center}

\newpage

%%%%%%%%%%%%%%%%%%%%%%%%%%%%%%%%%%%%%%%%%%%%%%%%%%%%%%%%%%%%%%%%%%%%%%%%%%%%%%%
% 10C CORE FRAMEWORK
%%%%%%%%%%%%%%%%%%%%%%%%%%%%%%%%%%%%%%%%%%%%%%%%%%%%%%%%%%%%%%%%%%%%%%%%%%%%%%%

\section*{Das 10C CORE Framework}
\addcontentsline{toc}{section}{Das 10C CORE Framework}

\begin{corebox}
BEATRIX basiert auf 9 Kernfragen, die jede Verhaltensentscheidung strukturieren:

\begin{center}
\begin{tabular}{clll}
\toprule
\textbf{\#} & \textbf{CORE} & \textbf{Frage} & \textbf{EBF-Kapitel} \\
\midrule
1 & WHO & Wer hat den Nutzen? & AAA \\
2 & WHAT & Was ist der Nutzen? (FEPSDE) & Kap. 10 \\
3 & HOW & Wie wirken Faktoren zusammen? ($\gamma$) & Kap. 5 \\
4 & WHEN & Wann zählt der Kontext? ($\Psi$) & Kap. 9 \\
5 & WHERE & Woher kommen die Parameter? & BBB \\
6 & AWARE & Was nimmt die Person wahr? & Kap. 11 \\
7 & READY & Wann ist sie handlungsbereit? & Kap. 12 \\
8 & STAGE & Wo steht sie in der Journey? & Kap. 13 \\
9 & HIERARCHY & Wie stratifizieren Entscheidungen? & HI \\
\bottomrule
\end{tabular}
\end{center}
\end{corebox}

\newpage

%%%%%%%%%%%%%%%%%%%%%%%%%%%%%%%%%%%%%%%%%%%%%%%%%%%%%%%%%%%%%%%%%%%%%%%%%%%%%%%
% FRAGE A: INNOVATION POTENTIAL
%%%%%%%%%%%%%%%%%%%%%%%%%%%%%%%%%%%%%%%%%%%%%%%%%%%%%%%%%%%%%%%%%%%%%%%%%%%%%%%

\section{Frage a) Innovation Potential}

\begin{keyinsight}
\textbf{BEATRIX ist eine Management-Innovation mit wissenschaftlicher Tiefe:} 8 formalisierte Kontext-Dimensionen, 15+ quantifizierte Wechselwirkungs-Parameter, 203+ dokumentierte Axiome.
\end{keyinsight}

\subsection{Das Problem, das wir lösen}

Stellen Sie sich vor: Ein Unternehmen will die Mitarbeiterzufriedenheit verbessern. Der Berater empfiehlt zehn Massnahmen -- aber welche davon wirken zusammen? Welche blockieren sich gegenseitig? Und funktioniert das, was in Skandinavien klappt, auch in der Schweiz?

\textbf{Heute basiert diese Einschätzung auf Erfahrung und Intuition.} BEATRIX macht sie systematisch und nachvollziehbar.

\subsection{Die drei Säulen der Innovation}

\subsubsection{Säule 1: Kontext ernst nehmen}

\begin{ebfbox}
\textbf{Kapitel 9} formalisiert den 8-dimensionalen Kontextvektor $\Psi$
\end{ebfbox}

Menschen verhalten sich nicht im luftleeren Raum. BEATRIX berücksichtigt \textbf{acht Kontext-Dimensionen} systematisch:

\begin{center}
\begin{tabular}{lll}
\toprule
\textbf{Dimension} & \textbf{Beschreibung} & \textbf{EBF-Code} \\
\midrule
Zeitlicher Kontext & Dringend oder langfristig? & $\Psi_T$ \\
Sozialer Kontext & Wer beobachtet? Wer urteilt? & $\Psi_S$ \\
Institutioneller Kontext & Welche Regeln? Welche Kultur? & $\Psi_I$ \\
Informationeller Kontext & Was weiss die Person? & $\Psi_D$ \\
Emotionaler Kontext & Gestresst? Entspannt? & $\Psi_E$ \\
Ressourcen-Kontext & Zeit, Geld, Energie? & $\Psi_E$ \\
Kultureller Kontext & Schweiz $\neq$ Schweden $\neq$ Singapur & $\Psi_C$ \\
Physischer Kontext & Wo findet Entscheidung statt? & $\Psi_P$ \\
\bottomrule
\end{tabular}
\end{center}

\textbf{Warum ist das neu?} Bisherige Ansätze behandeln Kontext als Störfaktor. Das EBF macht ihn zum zentralen Ordnungsprinzip mit 8 formalisierten Dimensionen.

\subsubsection{Säule 2: Wechselwirkungen sichtbar machen}

\begin{ebfbox}
\textbf{Kapitel 5} definiert die Komplementaritätsmatrix $C^*$ mit 15+ $\gamma$-Parametern
\end{ebfbox}

Wenn ein Unternehmen gleichzeitig Boni einführt und auf Teamarbeit setzt, können sich diese Massnahmen verstärken -- oder gegenseitig aufheben.

\begin{center}
\begin{tabular}{llll}
\toprule
\textbf{Kombination} & \textbf{Effekt} & \textbf{$\gamma$-Wert} & \textbf{Beispiel} \\
\midrule
Verstärkend & Wirken zusammen stärker & $\gamma > 0$ & Kommunikation + Feedback \\
Neutralisierend & Heben sich auf & $\gamma \approx 0$ & Unabhängige Massnahmen \\
Blockierend & Eine verhindert andere & $\gamma < 0$ & Boni + Teamziele \\
\bottomrule
\end{tabular}
\end{center}

\begin{formulabox}
\textbf{Die Komplementaritätsformel (Kapitel 5):}
\begin{equation}
E(\mathcal{I}_1 + \mathcal{I}_2) = E(\mathcal{I}_1) + E(\mathcal{I}_2) + \gamma_{12} \times \sqrt{E(\mathcal{I}_1) \times E(\mathcal{I}_2)}
\end{equation}
\end{formulabox}

\subsubsection{Säule 3: Transparente Entscheidungslogik}

\begin{ebfbox}
\textbf{Kapitel 11} enthält 104 Awareness-Axiome (AWX-A1 bis AWX-A90), die alle Empfehlungen begründen
\end{ebfbox}

BEATRIX ist kein ``Black Box''-System. Jede Empfehlung ist nachvollziehbar:

\begin{center}
\begin{tabular}{p{5.5cm}p{6.5cm}}
\toprule
\textbf{Traditionell} & \textbf{Mit BEATRIX} \\
\midrule
``Der Berater sagt, wir sollen X tun.'' & ``Die Analyse zeigt: In unserem Kontext wirkt X, weil Y. Evidenz aus Z.'' \\
Vertrauen in Person & Vertrauen in Methode \\
Schwer replizierbar & Auf 104+ Axiome rückführbar \\
\bottomrule
\end{tabular}
\end{center}

\subsection{Wissenschaftliche Fundierung: Das EBF}

Das Evidence-Based Framework integriert \textbf{50+ Jahre verhaltensökonomischer Forschung}:

\begin{center}
\begin{tabular}{llll}
\toprule
\textbf{Komponente} & \textbf{Kapitel} & \textbf{Inhalt} & \textbf{Quantifizierung} \\
\midrule
Komplementarität & 5 & Wechselwirkungen & 15+ $\gamma$-Parameter \\
Kontextmodellierung & 9 & Entscheidungskontext & 8 $\Psi$-Dimensionen \\
Welfare-Dimensionen & 10 & FEPSDE & 144 Utility-Komponenten \\
Awareness-Funktion & 11 & Wahrnehmung & 104 Axiome + 59 Korollare \\
Willingness-Funktion & 12 & Handlungsbereitschaft & 99 Axiome \\
Behavioral Change Journey & 13 & Verhaltensänderung & 6 Stufen, 152 Conditions \\
\bottomrule
\end{tabular}
\end{center}

\subsection{Was validiert der Forschungspartner?}

\textbf{Prof. Harald Gall} (Institut für Informatik):
\begin{itemize}[nosep]
    \item Prüft: Ist das System skalierbar?
    \item EBF-Bezug: Implementation von Kapitel 10 ($\Gamma$-Matrix 144×144)
\end{itemize}

\textbf{Prof. Johannes Luger} (Institut für Betriebswirtschaft):
\begin{itemize}[nosep]
    \item Prüft: Vertrauen Manager den Empfehlungen?
    \item EBF-Bezug: Forschung zu Kapitel 12 (WAX-Modell)
\end{itemize}

\begin{limitbox}
\textbf{Was BEATRIX nicht ist:}
\begin{itemize}[nosep]
    \item Kein Wahrsager -- keine individuellen Vorhersagen
    \item Kein Ersatz für Führung -- unterstützt, entscheidet nicht
    \item Kein Allheilmittel -- manche Probleme sind nicht analytisch lösbar
\end{itemize}
\end{limitbox}

\begin{summarybox}
\textbf{Innovation Potential:} BEATRIX formalisiert erstmals Kontext (8$\Psi$), Wechselwirkungen ($\gamma$-Matrix) und Entscheidungslogik (203+ Axiome) für die Beratungspraxis.
\end{summarybox}

\newpage

%%%%%%%%%%%%%%%%%%%%%%%%%%%%%%%%%%%%%%%%%%%%%%%%%%%%%%%%%%%%%%%%%%%%%%%%%%%%%%%
% FRAGE B: MARKET RELEVANCE
%%%%%%%%%%%%%%%%%%%%%%%%%%%%%%%%%%%%%%%%%%%%%%%%%%%%%%%%%%%%%%%%%%%%%%%%%%%%%%%

\section{Frage b) Market Relevance}

\begin{keyinsight}
\textbf{BEATRIX schliesst die ``Soft-Factor-Lücke'':} Für Finanzen gibt es SAP, für menschliches Verhalten fehlen systematische Werkzeuge. Das WAX $\gtrless$ $\theta$ Modell erklärt, warum.
\end{keyinsight}

\subsection{Die Soft-Factor-Lücke}

\begin{ebfbox}
\textbf{Kapitel 12} erklärt die ``Intentions-Verhaltens-Lücke'' formal
\end{ebfbox}

Moderne Unternehmen haben für fast alles Daten -- ausser für menschliches Verhalten:

\begin{center}
\begin{tabular}{lll}
\toprule
\textbf{Bereich} & \textbf{Werkzeuge} & \textbf{Was fehlt} \\
\midrule
Finanzen & SAP, Oracle & -- \\
Logistik & Supply-Chain-Systeme & -- \\
Marketing & CRM, Analytics & -- \\
\textbf{Mitarbeiterverhalten} & \textbf{Umfragen, Intuition} & \textbf{Systematik} \\
\textbf{Kundenentscheidungen} & \textbf{Fokusgruppen} & \textbf{Formale Modelle} \\
\bottomrule
\end{tabular}
\end{center}

\subsection{Die wissenschaftliche Erklärung}

\begin{formulabox}
\textbf{Die Handlungsgleichung (Kapitel 12):}
\begin{equation}
\text{Handlung} \Leftrightarrow WAX \geq \theta
\end{equation}
\begin{equation}
WAX = \varphi \times \min\{U^*\} + (1-\varphi) \times \sum U^*
\end{equation}
\end{formulabox}

\begin{center}
\begin{tabular}{lll}
\toprule
\textbf{Variable} & \textbf{Bedeutung} & \textbf{Praxis-Relevanz} \\
\midrule
$WAX$ & Willingness to Act & ``Wie bereit bin ich?'' \\
$\theta$ & Schwelle & ``Wie schwer ist es?'' \\
$\varphi$ & Denkstil & ``Denke ich verbunden oder getrennt?'' \\
$U^*$ & Wahrgenommener Nutzen & ``Was sehe ich?'' \\
\bottomrule
\end{tabular}
\end{center}

\textbf{Die Konsequenz:} Etwa 70\% aller Change-Projekte scheitern -- weil die WAX $\gtrless$ $\theta$ Dynamik nicht verstanden wird.

\subsection{Die 6 Stufen der Verhaltensänderung (BCJ)}

\begin{ebfbox}
\textbf{Kapitel 13} definiert 6 BCJ-Stufen für phasengerechte Interventionen
\end{ebfbox}

\begin{center}
\begin{tabular}{lllll}
\toprule
\textbf{Stufe} & \textbf{Name} & \textbf{Charakteristik} & \textbf{Intervention} \\
\midrule
$S_1$ & Unaware & Weiss nichts davon & Awareness schaffen \\
$S_2$ & Aware & Weiss, aber nicht relevant & Relevanz herstellen \\
$S_3$ & Considering & Überlegt & Benefits konkretisieren \\
$S_4$ & Intending & Will, aber tut nicht & Schwelle senken \\
$S_5$ & Acting & Handelt & Friktion minimieren \\
$S_6$ & Maintaining & Gewohnheit & Habit festigen \\
\bottomrule
\end{tabular}
\end{center}

\textbf{Praxis-Beispiel:}
\begin{itemize}[nosep]
    \item Anna in $S_3$ braucht konkrete Einsparungsberechnungen
    \item Thomas in $S_4$ braucht einen vereinfachten Antragsprozess
    \item Dieselbe Massnahme würde bei beiden unterschiedlich wirken
\end{itemize}

\subsection{Für wen ist BEATRIX gedacht?}

\begin{center}
\begin{tabular}{lll}
\toprule
\textbf{Wer?} & \textbf{Typische Frage} & \textbf{EBF-Anwendung} \\
\midrule
Beratungsunternehmen & ``Wie begründen wir Empfehlungen?'' & 104 Awareness-Axiome \\
Personalabteilungen & ``Warum funktioniert Strategie nicht?'' & 8$\Psi$ Kontext-Analyse \\
Öffentliche Verwaltung & ``Wie erreichen wir Bürger?'' & 6 BCJ-Stufen \\
Geschäftsführung & ``Welche Massnahmen wirken?'' & $\gamma$-Matrix Priorisierung \\
\bottomrule
\end{tabular}
\end{center}

\subsection{Der 5-Schritt-Prozess}

\begin{ebfbox}
\textbf{Kapitel 17} definiert die vollständige Interventionslogik
\end{ebfbox}

\begin{center}
\texttt{Kontext ($\Psi$) → 6D-Archetyp → 20-Field Schema → R-Score → GO/PILOT/NO-GO}
\end{center}

\begin{center}
\begin{tabular}{lll}
\toprule
\textbf{Schritt} & \textbf{Beschreibung} & \textbf{EBF-Komponente} \\
\midrule
1 & Kontext analysieren & 8 $\Psi$-Dimensionen \\
2 & Interventionstyp wählen & 8 Typen aus $W_{eff}$ \\
3 & Spezifikation ausfüllen & 20-Field Schema \\
4 & Relevanz berechnen & R-Score \\
5 & Entscheidung treffen & GO/PILOT/NO-GO \\
\bottomrule
\end{tabular}
\end{center}

\begin{summarybox}
\textbf{Market Relevance:} BEATRIX adressiert die ``Soft-Factor-Lücke'' mit WAX $\gtrless$ $\theta$ Modell (Kap. 12), 6 BCJ-Stufen (Kap. 13) und systematischem 5-Schritt-Prozess (Kap. 17).
\end{summarybox}

\newpage

%%%%%%%%%%%%%%%%%%%%%%%%%%%%%%%%%%%%%%%%%%%%%%%%%%%%%%%%%%%%%%%%%%%%%%%%%%%%%%%
% FRAGE C: FEASIBILITY & RISK REDUCTION
%%%%%%%%%%%%%%%%%%%%%%%%%%%%%%%%%%%%%%%%%%%%%%%%%%%%%%%%%%%%%%%%%%%%%%%%%%%%%%%

\section{Frage c) Feasibility \& Risk Reduction}

\begin{keyinsight}
\textbf{BEATRIX ist keine Idee auf dem Papier:} Das wissenschaftliche Fundament existiert -- 21 Kapitel, 203+ Axiome, 152 Boundary Conditions. Der Innovation Cheque prüft die praktische Anwendung.
\end{keyinsight}

\subsection{Das EBF als Fundament}

\begin{ebfbox}
21 Kapitel mit formaler Axiomatik sind bereits geschrieben
\end{ebfbox}

\begin{center}
\begin{tabular}{lll}
\toprule
\textbf{Komponente} & \textbf{Status} & \textbf{Quantifizierung} \\
\midrule
10C CORE Framework & ✓ Vollständig & 9 Kernfragen \\
Awareness-Axiome & ✓ Vollständig & 104 Axiome \\
Willingness-Axiome & ✓ Vollständig & 99 Axiome \\
BCJ-Modell & ✓ Vollständig & 6 Stufen, 152 Conditions \\
Interventionslogik & ✓ Vollständig & 8 Typen, 20-Field Schema \\
LLM-MC Methodik & ✓ Dokumentiert & Parameter-Schätzung \\
\bottomrule
\end{tabular}
\end{center}

\subsection{Was muss noch geprüft werden?}

\begin{center}
\begin{tabular}{llll}
\toprule
\textbf{Frage} & \textbf{EBF-Relevanz} & \textbf{Wer prüft?} \\
\midrule
Ist das System skalierbar? & Kap. 10: $\Gamma$-Matrix (144×144) & Prof. Gall \\
Vertrauen Manager den Empfehlungen? & Kap. 12: WAX-Modell & Prof. Luger \\
\bottomrule
\end{tabular}
\end{center}

\subsubsection{Technische Skalierbarkeit (Prof. Gall)}

\begin{center}
\begin{tabular}{lll}
\toprule
\textbf{Komponente} & \textbf{Komplexität} & \textbf{Prüfung nötig?} \\
\midrule
$\Gamma$-Matrix (144×144) & 20'736 Interaktionen & ✓ \\
AWX-Engine & 104 Axiome + 59 Korollare & ✓ \\
LLM-MC Pipeline & N=100+ Durchläufe & ✓ \\
20-Field Schema & Strukturierte Eingabe & ✓ \\
\bottomrule
\end{tabular}
\end{center}

\subsubsection{Nutzerakzeptanz (Prof. Luger)}

\begin{center}
\begin{tabular}{lll}
\toprule
\textbf{Forschungsfrage} & \textbf{EBF-Konzept} & \textbf{Testbar?} \\
\midrule
Wie beeinflusst $\varphi$ die Akzeptanz? & Kapitel 12 & ✓ \\
Wie verändert sich $\theta$ durch Human-AI? & Kapitel 12 & ✓ \\
Welche BCJ-Phasen sind kritisch? & Kapitel 13 & ✓ \\
\bottomrule
\end{tabular}
\end{center}

\subsection{Was kann mit CHF 15'000 erreicht werden?}

Der Innovation Cheque finanziert 125 Arbeitsstunden:

\begin{center}
\begin{tabular}{llll}
\toprule
\textbf{Arbeitspaket} & \textbf{Stunden} & \textbf{EBF-Bezug} \\
\midrule
Architektur-Analyse & 40 & Kap. 10 ($\Gamma$-Matrix) \\
Nutzer-Tests & 50 & Kap. 12 (WAX $\gtrless$ $\theta$) \\
Risikoanalyse & 20 & Kap. 17 (R-Score) \\
Abschlussbericht & 15 & Integration \\
\bottomrule
\end{tabular}
\end{center}

\subsection{Die LLM Monte Carlo Methodik}

\begin{ebfbox}
\textbf{Kapitel 10} dokumentiert die LLMMC-Methode für Parameter-Schätzung
\end{ebfbox}

\begin{center}
\begin{tabular}{ll}
\toprule
\textbf{Traditionell} & \textbf{Mit LLM-MC} \\
\midrule
Parameter via RCTs (Jahre) & Parameter via LLM-Simulation (Stunden) \\
Teure Primärforschung & Kostengünstige Schätzung \\
Einzelne Studien & N=100+ Durchläufe mit CI \\
\bottomrule
\end{tabular}
\end{center}

\textbf{Technische Details:}
\begin{itemize}[nosep]
    \item Mean estimation time: 0.3s per Parameter
    \item Output: 95\% Konfidenzintervalle
    \item Epistemic Tags: [EMP], [THR], [EXP]
\end{itemize}

\subsection{Risikomatrix}

\begin{center}
\begin{tabular}{lll}
\toprule
\textbf{Risiko} & \textbf{Wahrscheinlichkeit} & \textbf{EBF-Mitigation} \\
\midrule
Technische Skalierung & Mittel & Gall prüft $\Gamma$-Matrix \\
Manager misstrauen & Mittel & Luger testet WAX $\gtrless$ $\theta$ \\
Validierung fehlt & Niedrig & 203+ Axiome formalisiert \\
Vorhersagen falsch & Mittel & Axiom-Rückführbarkeit \\
\bottomrule
\end{tabular}
\end{center}

\begin{summarybox}
\textbf{Feasibility:} Das wissenschaftliche Fundament existiert (21 Kapitel, 203+ Axiome). Der Innovation Cheque prüft in 125h die praktische Skalierbarkeit und Akzeptanz.
\end{summarybox}

\newpage

%%%%%%%%%%%%%%%%%%%%%%%%%%%%%%%%%%%%%%%%%%%%%%%%%%%%%%%%%%%%%%%%%%%%%%%%%%%%%%%
% FRAGE D: RESEARCH PARTNER COLLABORATION
%%%%%%%%%%%%%%%%%%%%%%%%%%%%%%%%%%%%%%%%%%%%%%%%%%%%%%%%%%%%%%%%%%%%%%%%%%%%%%%

\section{Frage d) Research Partner Collaboration}

\begin{keyinsight}
\textbf{Die UZH bringt genau die Kompetenzen, die FehrAdvice fehlen:} Prof. Gall für technische Architektur (Kap. 10), Prof. Luger für Akzeptanzforschung (Kap. 12/13).
\end{keyinsight}

\subsection{Was FehrAdvice mitbringt}

\begin{center}
\begin{tabular}{ll}
\toprule
\textbf{Kompetenz} & \textbf{EBF-Bezug} \\
\midrule
Verhaltensökonomie & 50+ Jahre Fehr-Forschung integriert \\
Praxiswissen & 20+ Projekte als Pilotanwendungen \\
Ernst Fehr & CORE-Autor, Loss Aversion ($\lambda=5$) \\
\bottomrule
\end{tabular}
\end{center}

\subsection{Was FehrAdvice fehlt}

\begin{center}
\begin{tabular}{lll}
\toprule
\textbf{Lücke} & \textbf{UZH-Lösung} & \textbf{EBF-Kapitel} \\
\midrule
Technische Architektur & Prof. Gall & Kap. 10 \\
Akzeptanzforschung & Prof. Luger & Kap. 12 \\
Unabhängige Validierung & Beide & Kap. 17 \\
\bottomrule
\end{tabular}
\end{center}

\subsection{Prof. Gall: Technische Grundlagen}

\begin{ebfbox}
\textbf{EBF-Relevanz:} Implementation von Kapitel 10, 17
\end{ebfbox}

\begin{center}
\begin{tabular}{lll}
\toprule
\textbf{Beitrag} & \textbf{EBF-Komponente} & \textbf{Ergebnis} \\
\midrule
Architektur-Analyse & $\Gamma$-Matrix (144×144) & Skalierbarkeits-Report \\
Pipeline-Design & LLM-MC Methodik & Produktions-Architektur \\
20-Field Implementation & Kapitel 17 & Software-Spezifikation \\
\bottomrule
\end{tabular}
\end{center}

\subsection{Prof. Luger: Menschliche Faktoren}

\begin{ebfbox}
\textbf{EBF-Relevanz:} Forschung zu Kapitel 12, 13
\end{ebfbox}

\begin{center}
\begin{tabular}{lll}
\toprule
\textbf{Beitrag} & \textbf{EBF-Forschungsfrage} & \textbf{Ergebnis} \\
\midrule
Nutzer-Tests & ``Wie beeinflusst $\varphi$ die Akzeptanz?'' & Akzeptanz-Studie \\
Vertrauensforschung & ``Wann gilt WAX $>$ $\theta$ für KI?'' & Interventions-Design \\
Adoption-Phasen & ``Welche BCJ-Phase ist kritisch?'' & Change-Strategie \\
\bottomrule
\end{tabular}
\end{center}

\subsection{Wissenstransfer}

\begin{center}
\begin{tabular}{lll}
\toprule
\textbf{Bereich} & \textbf{Vorher} & \textbf{Nachher} \\
\midrule
Software-Architektur & Pragmatisch & EBF-konform (Kap. 10) \\
Akzeptanzforschung & Erfahrungsbasiert & WAX-Modell validiert (Kap. 12) \\
Change Management & Intuitiv & BCJ-basiert (Kap. 13) \\
\bottomrule
\end{tabular}
\end{center}

\begin{summarybox}
\textbf{Research Partner:} Gall validiert Kapitel 10 (Implementation), Luger validiert Kapitel 12/13 (Akzeptanz). Die UZH-Partnerschaft schliesst kritische Kompetenzlücken.
\end{summarybox}

\newpage

%%%%%%%%%%%%%%%%%%%%%%%%%%%%%%%%%%%%%%%%%%%%%%%%%%%%%%%%%%%%%%%%%%%%%%%%%%%%%%%
% FRAGE E: IMPACT FOR SWITZERLAND
%%%%%%%%%%%%%%%%%%%%%%%%%%%%%%%%%%%%%%%%%%%%%%%%%%%%%%%%%%%%%%%%%%%%%%%%%%%%%%%

\section{Frage e) Impact for Switzerland}

\begin{keyinsight}
\textbf{BEATRIX bringt Schweizer Verhaltensökonomie in den Markt:} 50+ Jahre Fehr-Forschung, 21 EBF-Kapitel, 203+ Axiome -- alles ``Made in Switzerland''.
\end{keyinsight}

\subsection{Eine Schweizer Lösung mit wissenschaftlicher Tiefe}

\begin{ebfbox}
\textbf{Kapitel 1} beschreibt Ernst Fehrs Vermächtnis
\end{ebfbox}

\begin{center}
\begin{tabular}{ll}
\toprule
\textbf{Aspekt} & \textbf{Bedeutung} \\
\midrule
Geistiges Eigentum & 21 Kapitel + 203 Axiome bleiben in der Schweiz \\
Fehrs Forschung & 50+ Jahre UZH-Verhaltensökonomie kommerzialisiert \\
Qualitätsstandard & ``Made in Switzerland'' für Beratungsmethodik \\
\bottomrule
\end{tabular}
\end{center}

\subsection{Wettbewerbsvorteil}

\begin{ebfbox}
\textbf{Kapitel 17} zeigt Differenzierung zu McKinsey/BCG
\end{ebfbox}

\begin{center}
\begin{tabular}{lll}
\toprule
\textbf{Aspekt} & \textbf{McKinsey/BCG} & \textbf{BEATRIX (EBF)} \\
\midrule
Methodik & Erfahrungsbasiert & 203+ Axiome \\
Vorhersagen & Qualitativ & Quantitativ (95\% CI) \\
Kontextanpassung & Manuell & Automatisch (8$\Psi$) \\
Kombinationseffekte & Ignoriert & $\gamma$-Matrix \\
Wissenschaftliche Basis & Schwach & 21 Kapitel \\
\bottomrule
\end{tabular}
\end{center}

\subsection{Die UZH→FehrAdvice Pipeline}

\begin{center}
\begin{tikzpicture}[node distance=1.2cm]
\node (uzh) {UZH-Forschung (Fehr)};
\node[below of=uzh] (ebf) {EBF Framework (21 Kapitel)};
\node[below of=ebf] (fa) {FehrAdvice Anwendung};
\node[below of=fa] (beatrix) {BEATRIX Tool};
\node[below of=beatrix] (export) {Schweizer Export};
\draw[->] (uzh) -- (ebf);
\draw[->] (ebf) -- (fa);
\draw[->] (fa) -- (beatrix);
\draw[->] (beatrix) -- (export);
\end{tikzpicture}
\end{center}

\textbf{Schweizer Standortvorteil:} Kurze Innovationspipeline von Wissenschaft zu Markt.

\subsection{Gesellschaftlicher Nutzen}

\begin{center}
\begin{tabular}{lll}
\toprule
\textbf{Bereich} & \textbf{Problem} & \textbf{EBF-Lösung} \\
\midrule
Gesundheit & Prävention wird wenig genutzt & BCJ-phasengerechte Gestaltung \\
Umwelt & Appelle verhallen & WAX $\gtrless$ $\theta$ Analyse \\
Finanzen & Viele sparen zu wenig & Default-Interventionen (Kap. 17) \\
Verwaltung & Angebote nicht genutzt & 8$\Psi$ Kontext-Optimierung \\
\bottomrule
\end{tabular}
\end{center}

\begin{summarybox}
\textbf{Impact Switzerland:} Die Fehr→UZH→FehrAdvice→BEATRIX Pipeline bringt 50+ Jahre Schweizer Verhaltensökonomie in den Markt -- mit 21 Kapiteln und 203+ Axiomen als Qualitätssiegel.
\end{summarybox}

\newpage

%%%%%%%%%%%%%%%%%%%%%%%%%%%%%%%%%%%%%%%%%%%%%%%%%%%%%%%%%%%%%%%%%%%%%%%%%%%%%%%
% GESAMTZUSAMMENFASSUNG
%%%%%%%%%%%%%%%%%%%%%%%%%%%%%%%%%%%%%%%%%%%%%%%%%%%%%%%%%%%%%%%%%%%%%%%%%%%%%%%

\section{Gesamtzusammenfassung}

\begin{center}
\begin{tabular}{p{1.5cm}p{6cm}p{4cm}}
\toprule
\textbf{Frage} & \textbf{Kernaussage} & \textbf{EBF-Kapitel} \\
\midrule
\textbf{a)} & Erstmals systematische Formalisierung: 8$\Psi$, $\gamma$-Matrix, 203+ Axiome & 5, 9, 10, 11 \\
\textbf{b)} & WAX $\gtrless$ $\theta$ erklärt die ``Soft-Factor-Lücke'' + BCJ für Targeting & 12, 13, 17 \\
\textbf{c)} & LLM-MC ermöglicht Parameter-Schätzung in 125h & 10, 17 \\
\textbf{d)} & Gall (Kap. 10) + Luger (Kap. 12) schliessen Lücken & 10, 12, 17 \\
\textbf{e)} & Fehr→UZH→FehrAdvice→BEATRIX Pipeline & 1, 17 \\
\bottomrule
\end{tabular}
\end{center}

\subsection*{Der strategische Kernsatz}

\begin{tcolorbox}[colback=innoblue!10!white, colframe=innoblue!75!black]
\centering
\textit{``Dieses Projekt entwickelt keine neue Software, sondern bringt 50 Jahre verhaltensökonomische Forschung -- formalisiert in 21 Kapiteln und 203+ Axiomen -- in die Beratungspraxis. Die Technologie ermöglicht diese Methodik, aber die Transformation der Entscheidungspraxis ist das eigentliche Produkt.''}
\end{tcolorbox}

\subsection*{EBF-Referenztabelle für den Antrag}

\begin{center}
\begin{tabular}{lllll}
\toprule
\textbf{Konzept} & \textbf{Kapitel} & \textbf{Axiome/Parameter} & \textbf{Für Frage} \\
\midrule
8$\Psi$ Kontextvektor & 9 & 8 Dimensionen & a, b \\
$\gamma$-Matrix & 5 & 15+ Parameter & a, c \\
FEPSDE & 10 & 144 Komponenten & a, d \\
AWX-Axiome & 11 & 104 Axiome & a, c \\
WAX $\gtrless$ $\theta$ & 12 & 99 Axiome & b, d \\
BCJ-Stufen & 13 & 6 Stufen, 152 Conditions & b \\
Interventionslogik & 17 & 8 Typen, 20-Field & b, c, d \\
R-Score & 17 & GO/PILOT/NO-GO & c \\
\bottomrule
\end{tabular}
\end{center}

\vfill

\begin{center}
\textcolor{innogray}{\rule{0.8\textwidth}{0.4pt}}\\[0.5em]
\textit{Dokument erstellt: 18. Januar 2026}\\
\textit{Version 3.0 -- Mit EBF-Kapitel-Integration}\\
\textit{Für: Innosuisse Innovation Cheque (SSBM)}
\end{center}

\end{document}
